\documentclass{article}
\usepackage[left=2cm, right=2cm, top=2cm, bottom=2cm]{geometry}
\usepackage{graphicx}
\usepackage{amsmath}
\usepackage{amssymb}
\usepackage{amsthm}
\usepackage{fancyhdr}
\usepackage{verbatim}
\usepackage{listings}
\usepackage{xcolor}
\usepackage{pdfpages}
\usepackage{cancel}
\usepackage{physics}
\usepackage{esint}

\lstset{
    language=C++,
    basicstyle=\ttfamily\footnotesize,
    keywordstyle=\color{blue},
    commentstyle=\color{green},
    stringstyle=\color{red},
    numbers=left,
    numberstyle=\tiny\color{gray},
    frame=single,
    breaklines=true,
    captionpos=b,
}


\title{PHYS 487: Lecture \# 21}
\author{Cliff Sun}

\newtheorem{theorem}{Theorem}[section]
\newtheorem{lemma}[theorem]{Lemma}
\newtheorem{definition}[theorem]{Definition}
\newtheorem{conjecture}[theorem]{Conjecture}
\newtheorem{proposition}[theorem]{Proposition}
\newtheorem{corollary}[theorem]{Corollary}
\newtheorem{one minute paper}[theorem]{One Minute Paper}

\pagestyle{fancy}
\lhead{\textbf{\thepage}\ \ \nouppercase{\rightmark}}
\chead{PHYS 487: Lecture \# 21}
\rhead{Cliff Sun}

\begin{document}

\maketitle

\section*{Lecture Span}
\begin{enumerate}
    \item More quantum stuff
\end{enumerate}

\section*{Density Matrix}

Density matrix is defined as 

\begin{equation}
    \rho: \;\; \ket{\psi}\bra{\psi}
\end{equation}

Note that $\rho_{ii}$ are all real values and 

\begin{equation}
    \sum_i \rho_{ii} =1 \iff \Tr(\rho) = 1
\end{equation}

Moreover, 

\begin{equation}
    \rho_{ij} = \rho_{ji}^{*}
\end{equation}

\section*{Imperfect Knowledge}

If we want to move the system into a state $\ket{+}$ but there is some variation from the desired state. Then how do we describe this? Use the density matrix 

\begin{equation}
    \rho = \frac{1}{N}\sum_{\chi}\ket{\psi_{\chi}}\bra{\psi_{\chi}}
\end{equation}

for mixed states. That is $\ket{\psi_{\chi}} \neq \ket{\psi_{i}}$ for all $i, \chi$. Here, we don't know $\ket{\psi_{\chi}}$ exactly. Then we can represent everything with an average vector. Here, 

\begin{enumerate}
    \item totally mixed: average is $0$
    \item totally pure: average is 1
    \item a bit of both: average $\in (0,1)$
\end{enumerate}

Pure states have a basis where 

\begin{equation}
    \rho^2 = \rho
\end{equation}

This is also called an \underline{Idempotent matrix}. Also in pure states, 

\begin{equation}
    \Tr[\rho^2] = \Tr[\rho] =1 
\end{equation}

But for mixed states 

\begin{equation}
    \Tr[\rho^2] < \Tr[\rho] < 1
\end{equation}

Then we can define a linear entropy (mixing) as 

\begin{equation}
    1 - \Tr[\rho^2]
\end{equation}

We can also define measurements, 

\begin{equation}
    \Pr[\chi] = | \langle \chi | \psi \rangle |^2 = \langle \chi | \rho | \psi \rangle
\end{equation}

And an expectation value 

\begin{equation}
    \langle A \rangle = \Tr[\rho A]
\end{equation}

And a time evolution on $\rho$ as 

\begin{equation}
    \rho = -\frac{i}{\chi}[\rho, H]
\end{equation}


\end{document}